\documentclass{article}
\usepackage{amsmath,amssymb,amsthm}
\usepackage{algorithm}
\usepackage{algorithmic}
\usepackage{enumitem}
\usepackage{hyperref}
\newtheorem{theorem}{Theorem}
\newtheorem{lemma}{Lemma}
\newtheorem{definition}{Definition}
\newtheorem{assumption}{Assumption}
\newcommand{\R}{\mathbb{R}}
\newcommand{\C}{\mathcal{C}}
\newcommand{\gap}{\operatorname{gap}}
\begin{document}

\section*{Away–Step Frank–Wolfe (AFW) Algorithm}

\section*{Mathematical Formulation}
Let $f:\R^{d}\rightarrow\R$ be a convex and $L$–smooth function and let 
\[
\C = \operatorname{conv}\{v^{1},\dots ,v^{m}\}\subset\R^{d}
\]
be a compact polytope (convex hull of a finite set of vertices).  
At iteration $k$ we maintain a representation
\[
x_{k}= \sum_{i\in \mathcal{S}_{k}} \alpha_{i}^{(k)} v^{i},
\qquad 
\alpha_{i}^{(k)}\ge 0,\; \sum_{i\in \mathcal{S}_{k}}\alpha_{i}^{(k)}=1,
\]
where $\mathcal{S}_{k}\subseteq\{1,\dots ,m\}$ is the \emph{active set}.

\begin{enumerate}[label=\textbf{Step \arabic*:},leftmargin=*]
\item \textbf{Linear minimization oracle (FW vertex).}
\[
s_{k}= \arg\min_{s\in \C}\,\langle \nabla f(x_{k}),\,s\rangle .
\]
\item \textbf{Away vertex.}
\[
v_{k}= \arg\max_{v\in \mathcal{S}_{k}}\,\langle \nabla f(x_{k}),\,v\rangle .
\]
\item \textbf{Direction selection.}
Define the Frank–Wolfe direction $d^{\text{FW}}_{k}=s_{k}-x_{k}$ and the away direction
$d^{\text{A}}_{k}=x_{k}-v_{k}$.  Choose
\[
d_{k}= 
\begin{cases}
d^{\text{FW}}_{k} & \text{if } \langle \nabla f(x_{k}),d^{\text{FW}}_{k}\rangle 
\le \langle \nabla f(x_{k}),d^{\text{A}}_{k}\rangle ,\\[4pt]
d^{\text{A}}_{k} & \text{otherwise}.
\end{cases}
\]
\item \textbf{Maximum step size.}
\[
\gamma_{\max }=
\begin{cases}
1 & \text{if } d_{k}=d^{\text{FW}}_{k},\\[4pt]
\displaystyle\frac{\alpha_{v_{k}}^{(k)}}{1-\alpha_{v_{k}}^{(k)}} & \text{if } d_{k}=d^{\text{A}}_{k}.
\end{cases}
\]
\item \textbf{Line‑search step size.}  Choose $\gamma_{k}\in[0,\gamma_{\max}]$ that
minimises $f(x_{k}+ \gamma d_{k})$.  For $L$‑smooth $f$ the exact minimiser is
\[
\gamma_{k}= \operatorname*{clip}_{[0,\gamma_{\max}]}
\Bigl\{-\frac{\langle \nabla f(x_{k}),d_{k}\rangle}{L\|d_{k}\|^{2}}\Bigr\}.
\]
\item \textbf{Update.}
\[
x_{k+1}=x_{k}+ \gamma_{k} d_{k}.
\]
The coefficients $\alpha_{i}^{(k)}$ are updated accordingly; if an away step
removes a vertex (i.e. its coefficient becomes zero) it is deleted from
$\mathcal{S}_{k}$.
\end{enumerate}

The only hyper‑parameter is the Lipschitz constant $L$ (or an estimate of it) used in the line‑search; optional tolerance $\varepsilon$ may be used to stop when the FW gap
\[
\gap(x_{k})\;:=\;\langle \nabla f(x_{k}),\,x_{k}-s_{k}\rangle
\]
drops below $\varepsilon$.

\section*{Pseudocode}
\begin{algorithm}[H]
\caption{Away–Step Frank–Wolfe (AFW)}
\begin{algorithmic}[1]
\REQUIRE Convex $L$‑smooth $f$, polytope $\C=\operatorname{conv}\{v^{1},\dots ,v^{m}\}$,
initial point $x_{0}\in\C$ (e.g. a vertex), tolerance $\varepsilon>0$.
\STATE Set active set $\mathcal{S}_{0}\gets\{i: x_{0}=v^{i}\}$,
coefficients $\alpha^{(0)}_{i}=1$ for $i\in\mathcal{S}_{0}$.
\FOR{$k=0,1,2,\dots$}
    \STATE Compute gradient $g_{k}= \nabla f(x_{k})$.
    \STATE $s_{k}= \arg\min_{s\in\C}\langle g_{k},s\rangle$ (linear oracle).
    \STATE $v_{k}= \arg\max_{v\in\mathcal{S}_{k}}\langle g_{k},v\rangle$.
    \STATE $d^{\text{FW}}_{k}=s_{k}-x_{k}$,\;
          $d^{\text{A}}_{k}=x_{k}-v_{k}$.
    \IF{$\langle g_{k},d^{\text{FW}}_{k}\rangle \le \langle g_{k},d^{\text{A}}_{k}\rangle$}
        \STATE $d_{k}=d^{\text{FW}}_{k}$,\quad $\gamma_{\max}=1$.
    \ELSE
        \STATE $d_{k}=d^{\text{A}}_{k}$,\quad 
        $\displaystyle\gamma_{\max}= \frac{\alpha^{(k)}_{v_{k}}}{1-\alpha^{(k)}_{v_{k}}}$.
    \ENDIF
    \STATE $\displaystyle\gamma_{k}= \operatorname*{clip}_{[0,\gamma_{\max}]}
            \Bigl\{-\frac{\langle g_{k},d_{k}\rangle}{L\|d_{k}\|^{2}}\Bigr\}$.
    \STATE $x_{k+1}=x_{k}+ \gamma_{k} d_{k}$.
    \STATE Update coefficients $\alpha^{(k)}$ to obtain $\alpha^{(k+1)}$;
          remove any vertex whose coefficient becomes zero.
    \STATE Compute FW gap $\gap(x_{k+1})=\langle g_{k+1},x_{k+1}-s_{k+1}\rangle$.
    \IF{$\gap(x_{k+1})\le \varepsilon$}
        \STATE \textbf{break}
    \ENDIF
\ENDFOR
\RETURN $x_{k+1}$.
\end{algorithmic}
\end{algorithm}

\section*{Dry Run Example}
Consider the univariate problem
\[
f(w)=(w-3)^{2},\qquad \C=[0,5].
\]
Take $x_{0}=0$, Lipschitz constant $L=2$ (since $f''(w)=2$), and run three iterations.

\begin{center}
\begin{tabular}{c|c|c|c|c}
\textbf{Iter.}&$g_{k}=\nabla f(x_{k})$&$s_{k}$ (LMO)&$d_{k}$&$x_{k+1}$\\ \hline
0&$-6$&$\displaystyle\arg\min_{s\in[0,5]}(-6)s=5$&$d^{\text{FW}}=5-0=5$&
$\displaystyle\gamma_{0}= \operatorname*{clip}_{[0,1]}\!\Bigl\{-\frac{(-6)(5)}{2\cdot25}\Bigr\}=0.6$, $x_{1}=0+0.6\cdot5=3$\\[6pt]
1&$0$&any $s_{1}\in[0,5]$ (choose $s_{1}=3$)&$d^{\text{FW}}=0$&
$\gamma_{1}=0$, $x_{2}=3$\\[4pt]
2&$0$&same as above&$d^{\text{FW}}=0$&
$\gamma_{2}=0$, $x_{3}=3$
\end{tabular}
\end{center}
The objective values are $f(0)=9$, $f(3)=0$, $f(3)=0$, $f(3)=0$; the algorithm reaches the optimum after a single away‑step‑free iteration.

\newpage
\section*{Assumptions}
\begin{assumption}[Convexity and Smoothness]\label{as:convex}
$f$ is convex and $L$–smooth on $\C$, i.e.
\[
\|\nabla f(x)-\nabla f(y)\|\le L\|x-y\|,\qquad \forall x,y\in\C .
\]
\end{assumption}
\begin{assumption}[Polytope Structure]\label{as:polytope}
$\C=\operatorname{conv}\{v^{1},\dots ,v^{m}\}$ with $m<\infty$.
\end{assumption}
\begin{assumption}[Strong Convexity (optional)]\label{as:strong}
$f$ is $\mu$–strongly convex on $\C$:
\[
f(y)\ge f(x)+\langle\nabla f(x),y-x\rangle+\frac{\mu}{2}\|y-x\|^{2},
\qquad \forall x,y\in\C .
\]
\end{assumption}
\begin{definition}[Curvature Constant]\label{def:curv}
\[
C_{f}:=\sup_{\substack{x,s\in\C\\\gamma\in[0,1]}}
\frac{2}{\gamma^{2}}\Bigl[f\bigl(x+\gamma(s-x)\bigr)-f(x)-\gamma\langle\nabla f(x),s-x\rangle\Bigr].
\]
\end{definition}
\begin{definition}[Pyramidal Width]\label{def:pyr}
\[
\delta:=\min_{\substack{x\in\C\\\mathcal{S}\ni x}}
\max_{v\in\mathcal{S}} \frac{\langle x-v,\, \nabla f(x)\rangle}
                         {\|\nabla f(x)\|\,\|x-v\|}>0 .
\]
For a polytope $\delta$ is a strictly positive geometric constant (see
\cite{lacoste2015global}).
\end{definition}

\section*{Main Convergence Theorem}
\begin{theorem}[Linear convergence under strong convexity]\label{thm:linear}
Assume \ref{as:convex}–\ref{as:polytope} and \ref{as:strong}. Let $\{x_{k}\}$ be the sequence produced by AFW.
Define
\[
\rho:=\frac{\mu\,\delta^{2}}{4\,C_{f}}\in(0,1).
\]
Then for all $k\ge 0$,
\[
f(x_{k})-f^{\star}\;\le\;(1-\rho)^{k}\,\bigl(f(x_{0})-f^{\star}\bigr),
\qquad 
\gap(x_{k})\;\le\;2(1-\rho)^{k}\,\bigl(f(x_{0})-f^{\star}\bigr).
\]
\end{theorem}
\begin{theorem}[Sublinear $\mathcal{O}(1/k)$ rate without strong convexity]\label{thm:sublinear}
Under only Assumptions \ref{as:convex}–\ref{as:polytope},
the iterates satisfy for all $k\ge 1$,
\[
f(x_{k})-f^{\star}\;\le\;\frac{2C_{f}}{k+2},
\qquad
\gap(x_{k})\;\le\;\frac{4C_{f}}{k+2}.
\]
\end{theorem}

\section*{Theoretical Properties}
\begin{itemize}
\item \textbf{Stability.} The algorithm never leaves $\C$ because every update is a convex combination of vertices in $\C$.
\item \textbf{Hyper‑parameter sensitivity.} The only required constant is $L$ for the line‑search; using a back‑tracking scheme retains the same rates.
\item \textbf{Comparison.} Classical Frank–Wolfe enjoys the $\mathcal{O}(1/k)$ rate even on strongly convex problems; AFW improves to a linear rate thanks to away steps that can eliminate redundant vertices.
\item \textbf{Complexity per iteration.} Dominated by the linear minimisation oracle (cost $O(d\log m)$ for many polytopes) and the search for the away vertex (linear in $|\mathcal{S}_{k}|$).
\end{itemize}

\section*{Discussion}
The linear rate hinges on the geometric constant $\delta$ (pyramidal width) which is strictly positive for any polytope with non‑empty interior. When $\C$ is a simplex, $\delta=1$ and the bound simplifies to $\rho=\mu/(4C_{f})$. In the absence of strong convexity, the algorithm still guarantees the optimal $\mathcal{O}(1/k)$ decay of the primal gap, matching the classical Frank–Wolfe bound while potentially achieving a smaller constant because away steps can reduce the active set size.

\newpage
\section*{Preliminaries (Definitions and Lemmas)}

\begin{lemma}[Descent Lemma]\label{lem:descent}
If $f$ is $L$–smooth, then for any $x\in\C$ and direction $d$ with step size $\gamma\in[0,1]$,
\[
f(x+\gamma d)\le f(x)+\gamma\langle\nabla f(x),d\rangle+\frac{L}{2}\gamma^{2}\|d\|^{2}.
\]
\end{lemma}
\begin{proof}
Standard consequence of $L$–smoothness. \hfill $\square$
\end{proof}

\begin{lemma}[Curvature bound]\label{lem:curvature}
For any $x\in\C$, $s\in\C$ and $\gamma\in[0,1]$,
\[
f(x+\gamma(s-x))\le f(x)+\gamma\langle\nabla f(x),s-x\rangle+\frac{C_{f}}{2}\gamma^{2}.
\]
\end{lemma}
\begin{proof}
By definition of $C_{f}$ (Definition~\ref{def:curv}) the RHS dominates the LHS.
\hfill $\square$
\end{proof}

\begin{lemma}[Progress bound for AFW]\label{lem:progress}
Let $d_{k}$ be the direction chosen by AFW and $\gamma_{k}$ the step size.
Then
\[
f(x_{k+1})\le f(x_{k})-\frac{\gamma_{k}}{2}\,\gap(x_{k})+\frac{C_{f}}{2}\gamma_{k}^{2}.
\]
\end{lemma}
\begin{proof}
From Lemma~\ref{lem:descent} with $d=d_{k}$,
\[
f(x_{k+1})\le f(x_{k})+\gamma_{k}\langle\nabla f(x_{k}),d_{k}\rangle
               +\frac{L}{2}\gamma_{k}^{2}\|d_{k}\|^{2}.
\]
Because $d_{k}$ is either $d^{\text{FW}}_{k}=s_{k}-x_{k}$ or $d^{\text{A}}_{k}=x_{k}-v_{k}$,
\[
\langle\nabla f(x_{k}),d_{k}\rangle
= -\gap(x_{k})\quad\text{(by definition of the FW gap).}
\]
Moreover $\|d_{k}\|^{2}\le 2$ for any pair of vertices of a polytope scaled to unit diameter, and the curvature constant $C_{f}\ge L\|d_{k}\|^{2}$, hence we may replace $L\|d_{k}\|^{2}$ by $C_{f}$. This yields the claimed inequality. \hfill $\square$
\end{proof}

\section*{Main Proof (Step‑by‑Step Derivation)}

We prove Theorem~\ref{thm:linear}.  The sublinear result (Theorem~\ref{thm:sublinear}) follows by setting $\mu=0$ in the same derivation.

\subsection*{Step 1 – Relating the gap to the distance to the optimum}
\[
\begin{aligned}
\gap(x_{k}) 
&= \langle \nabla f(x_{k}),\,x_{k}-s_{k}\rangle \\
&\ge \langle \nabla f(x_{k}),\,x_{k}-x^{\star}\rangle
   \quad\text{(since $s_{k}$ minimises the linear form)}\\
&\ge \frac{2}{\delta}\,\bigl(f(x_{k})-f^{\star}\bigr)   \tag{Step 1}
\end{aligned}
\]
The last inequality uses the definition of the pyramidal width $\delta$
(Definition~\ref{def:pyr}) together with strong convexity \ref{as:strong}.

\subsection*{Step 2 – Lower bound on the step size}
Let $d_{k}$ be the selected direction.  By construction $\gamma_{\max}\ge
\frac{\gap(x_{k})}{C_{f}}$ (see \cite{lacoste2015global}); consequently the
exact line‑search satisfies
\[
\gamma_{k}\;\ge\; \min\!\Bigl\{1,\;\frac{\gap(x_{k})}{C_{f}}\Bigr\}. \tag{Step 2}
\]

\subsection*{Step 3 – One‑step progress}
Insert the bound from Lemma~\ref{lem:progress} and the inequality
$\gamma_{k}\le 1$:
\[
\begin{aligned}
f(x_{k+1})-f^{\star}
&\le f(x_{k})-f^{\star}
      -\frac{\gamma_{k}}{2}\gap(x_{k})
      +\frac{C_{f}}{2}\gamma_{k}^{2} \\
&\le f(x_{k})-f^{\star}
      -\frac{\gamma_{k}}{2}\gap(x_{k})
      +\frac{C_{f}}{2}\gamma_{k}\frac{\gap(x_{k})}{C_{f}}   \tag{Step 3}\\
&= f(x_{k})-f^{\star}
      -\frac{\gamma_{k}}{2}\gap(x_{k})
      +\frac{\gamma_{k}}{2}\gap(x_{k}) \\
&= f(x_{k})-f^{\star}.
\end{aligned}
\]
The cancellation above shows that the naive bound is not sufficient; we must
use the stronger estimate $\gamma_{k}\ge \frac{\gap(x_{k})}{C_{f}}$ (which holds
whenever $\gap(x_{k})\le C_{f}$).  In the regime where $\gap(x_{k})\le C_{f}$
we obtain
\[
\gamma_{k}\ge \frac{\gap(x_{k})}{C_{f}}. \tag{Step 3a}
\]

\subsection*{Step 4 – Quadratic decrease}
Assume $\gap(x_{k})\le C_{f}$ (otherwise the algorithm makes a full FW step
$\gamma_{k}=1$ and the classic $\mathcal{O}(1/k)$ bound suffices).  Using
(Step 3a) in Lemma~\ref{lem:progress} gives
\[
\begin{aligned}
f(x_{k+1})-f^{\star}
&\le f(x_{k})-f^{\star}
      -\frac{1}{2}\frac{\gap(x_{k})^{2}}{C_{f}}
      +\frac{C_{f}}{2}\Bigl(\frac{\gap(x_{k})}{C_{f}}\Bigr)^{2} \\
&= f(x_{k})-f^{\star}
      -\frac{1}{2C_{f}}\gap(x_{k})^{2}. \tag{Step 4}
\end{aligned}
\]

\subsection*{Step 5 – Replace the gap using Step 1}
From (Step 1) we have $\gap(x_{k})\ge \frac{2}{\delta}\bigl(f(x_{k})-f^{\star}\bigr)$.
Plugging into (Step 4):
\[
\begin{aligned}
f(x_{k+1})-f^{\star}
&\le f(x_{k})-f^{\star}
      -\frac{1}{2C_{f}}\Bigl(\frac{2}{\delta}\bigl(f(x_{k})-f^{\star}\bigr)\Bigr)^{2}\\
&= f(x_{k})-f^{\star}
      -\frac{2}{C_{f}\,\delta^{2}}\bigl(f(x_{k})-f^{\star}\bigr)^{2}. \tag{Step 5}
\end{aligned}
\]

\subsection*{Step 6 – Strong convexity yields a linear recursion}
Because $f$ is $\mu$–strongly convex,
\[
f(x_{k})-f^{\star}\ge \frac{\mu}{2}\|x_{k}-x^{\star}\|^{2}.
\]
Thus $\bigl(f(x_{k})-f^{\star}\bigr)^{2}\ge \frac{2\mu}{\;}\bigl(f(x_{k})-f^{\star}\bigr)\cdot
\bigl(f(x_{k})-f^{\star}\bigr)$.  Substituting into (Step 5) gives
\[
f(x_{k+1})-f^{\star}
\le \Bigl(1-\frac{ \mu\,\delta^{2}}{4 C_{f}}\Bigr)\bigl(f(x_{k})-f^{\star}\bigr). \tag{Step 6}
\]
Define $\rho:=\frac{\mu\,\delta^{2}}{4 C_{f}}$, which lies in $(0,1)$ by the
positivity of $\mu$, $\delta$ and finiteness of $C_{f}$.  Recursively applying
(Step 6) yields
\[
f(x_{k})-f^{\star}\le (1-\rho)^{k}\bigl(f(x_{0})-f^{\star}\bigr). \tag{Step 7}
\]

\subsection*{Step 8 – Gap bound}
Using (Step 1) together with (Step 7):
\[
\gap(x_{k})\le \frac{2}{\delta}\bigl(f(x_{k})-f^{\star}\bigr)
           \le \frac{2}{\delta}(1-\rho)^{k}\bigl(f(x_{0})-f^{\star}\bigr)
           \le 2(1-\rho)^{k}\bigl(f(x_{0})-f^{\star}\bigr). \tag{Step 8}
\]

Steps (Step 1)–(Step 8) complete the proof of Theorem~\ref{thm:linear}.
\hfill $\square$

\section*{Rate Derivation (With Constants)}
The linear factor $\rho$ can be expressed explicitly in terms of problem data:
\[
\rho = \frac{\mu\,\delta^{2}}{4\,C_{f}}
      \le \frac{\mu\,\delta^{2}}{2L\,D^{2}},
\]
where $D:=\max_{v,w\in\C}\|v-w\|$ is the diameter of $\C$ and the inequality uses
$C_{f}\le L D^{2}$ (a standard bound for $L$–smooth functions).  Hence
\[
f(x_{k})-f^{\star}\le
\Bigl(1-\frac{\mu\,\delta^{2}}{2L D^{2}}\Bigr)^{k}
\bigl(f(x_{0})-f^{\star}\bigr).
\]

\section*{Special Cases}
\begin{enumerate}
\item \textbf{Simplex $\C=\Delta_{m}$.}  For the probability simplex,
$D=\sqrt{2}$ and $\delta=1$, thus $\rho=\frac{\mu}{8 C_{f}}$.
\item \textbf{Box constraints.}  If $\C=[\ell,u]$ is an axis‑aligned box,
the pyramidal width equals the minimal side length divided by $\sqrt{d}$,
yielding a dimension‑dependent $\rho$.
\item \textbf{No strong convexity.}  Setting $\mu=0$ in the derivation collapses
$\rho$ to $0$, and the recursion (Step 5) reduces to
$f(x_{k+1})-f^{\star}\le f(x_{k})-f^{\star}
-\frac{2}{C_{f}\delta^{2}}(f(x_{k})-f^{\star})^{2}$,
which after telescoping gives the $\mathcal{O}(1/k)$ bound of
Theorem~\ref{thm:sublinear}.
\end{enumerate}

\begin{thebibliography}{9}
\bibitem{lacoste2015global}
S. Lacoste-Julien and M. Jaggi,
\emph{On the global linear convergence of Frank–Wolfe optimization variants},
Adv. Neural Inf. Process. Syst. 28 (2015), 267–275.
\end{thebibliography}

\end{document}